\documentclass{szb-solution}

\area{Kalkulus}
\title{Differenciálás I}
\subject{Matematika G1}
\subjectCode{BMETE94BG01}
\date{Utoljára frissítve: \today}
\docno{7}

\begin{document}
\maketitle

\subsection{A hatványfüggvény deriváltja a definícióból}

Legyen $f(x)=x^n$, ahol $n\in\mathbb N$. A differenciálhányados definíciója
szerint
$$
  f'(x_0)=\lim_{x\to x_0}\frac{x^n-x_0^n}{x-x_0}
  \text.
$$
Felhasználjuk a hatványok különbségének szorzattá alakítását:
$$
  x^n-x_0^n=(x-x_0)
  \left(x^{n-1}+x^{n-2}x_0+\dots+xx_0^{n-2}+x_0^{n-1}\right)
  \text.
$$
Ha $x\neq x_0$, egyszerűsíthetünk $x-x_0$-lal, majd határértéket vehetünk:
\begin{align*}
  f'(x_0)
  &=\lim_{x\to x_0}
    \left(x^{n-1}+x^{n-2}x_0+\dots+x_0^{n-1}\right)\\
  &=\underbrace{x_0^{n-1}+x_0^{n-1}+\dots+x_0^{n-1}}_{n\text{ darab}}
  =n x_0^{n-1}
  \text.
\end{align*}
Tehát
$$
  \boxed{(x^n)'=nx^{n-1}}
  \text.
$$

\subsection{Szakaszonként megadott függvények differenciálhatósága}

\begin{enumerate}[a)]
  \item A függvény mindkét oldalról nullához tart, és $f(0)=0$, tehát
        folytonos a nullában. Az egyoldali differenciálhányadosok:
        $$
          f'_-(0)=\lim_{x\to0^-}\frac{\sin^2x}{x}
          =\lim_{x\to0^-}\frac{\sin x}{x}\sin x=0,
        $$
        $$
          f'_+(0)=\lim_{x\to0^+}\frac{x^2}{x}
          =\lim_{x\to0^+}x=0
          \text.
        $$
        A két érték megegyezik, ezért
        $$
          \boxed{f\text{ differenciálható }0\text{-ban, és }f'(0)=0.}
        $$
        A derivált teljes alakja
        $$
          f'(x)=
          \begin{cases}
            \sin2x,&x<0,\\
            0,&x=0,\\
            2x,&x>0.
          \end{cases}
        $$

  \item A jobb oldali határérték $0$, a bal oldali viszont
        $$
          \lim_{x\to0^-}(x^3+x+1)=1
          \text.
        $$
        Ezek nem egyeznek meg, ezért a függvény a nullában nem folytonos.
        A differenciálhatóság folytonosságot von maga után, így
        $$
          \boxed{f\text{ nem differenciálható }0\text{-ban}.}
        $$
\end{enumerate}

\subsection{Példák folytonosságra és differenciálhatóságra}

\begin{enumerate}[a)]
  \item Az $f(x)=|x-1|$ függvény mindenütt folytonos, de $x=1$-ben a bal
        oldali derivált $-1$, a jobb oldali derivált $1$, tehát ott nem
        differenciálható.

  \item Ilyen függvény nem létezik. Ha egy függvény differenciálható egy
        pontban, akkor abban a pontban szükségképpen folytonos is.

  \item Például $f(x)=x^2$. Ekkor $f'(x)=2x$, így mind $f$, mind $f'$ az egész
        valós számegyenesen folytonos.

  \item Legyen
        $$
          f(x)=
          \begin{cases}
            x^2\sin(1/x),&x\neq0,\\
            0,&x=0.
          \end{cases}
        $$
        A nullában
        $$
          f'(0)=\lim_{x\to0}\frac{x^2\sin(1/x)}x
          =\lim_{x\to0}x\sin(1/x)=0
          \text,
        $$
        tehát $f$ mindenütt differenciálható. Ha $x\neq0$, akkor
        $$
          f'(x)=2x\sin(1/x)-\cos(1/x)
          \text.
        $$
        Ez $x\to0$ esetén nem tart nullához, ezért $f'$ nem folytonos a
        nullában.
\end{enumerate}

\subsection{Differenciálhatóság és közeli töréspontok}

Legyen
$$
  f(x)=
  \begin{cases}
    x^2|\sin(1/x)|,&x\neq0,\\
    0,&x=0.
  \end{cases}
$$
A nullában
$$
  \left|\frac{f(x)-f(0)}{x}\right|
  =|x|\,|\sin(1/x)|\leq|x|\longrightarrow0
  \text,
$$
tehát $f'(0)=0$: a függvény differenciálható a nullában.

Ugyanakkor $\sin(1/x)=0$, ha
$$
  x=x_k=\frac1{k\pi},
  \qquad k\in\mathbb Z\setminus\{0\}
  \text.
$$
Ezekben a pontokban $\sin(1/x)$ egyszeres zérushelyű, ezért az abszolútérték
bal és jobb oldali meredeksége ellentétes előjelű. Az $x^2$ pozitív tényező ezt
a törést nem szünteti meg, így $f$ egyik $x_k$ pontban sem differenciálható.
Mivel $x_k\to0$, a nulla bármilyen kis környezetében található ilyen pont.

\subsection{Deriválási gyakorlatok}

\begin{enumerate}[a)]
  \item A $\sin^2x+\cos^2x=1$ tag deriváltja nulla, ezért
        $$
          \boxed{
          f'(x)=5(6x^7+7x^4+2x^2)^4(42x^6+28x^3+4x)}
          \text.
        $$

  \item A szorzat- és hatványszabályt alkalmazva
        $$
          \boxed{
          g'(x)=\frac{e^x}{x}+e^x\ln x
          +2x\cot x-\frac{x^2}{\sin^2x}-\frac13x^{-4/3}}
          \text.
        $$

  \item Jelölje $C=\operatorname{arctanh}\pi$ a nevezőben szereplő
        konstans értéket. Az eredeti kifejezést formálisan differenciálva,
        logaritmikus deriválással tömör alakot kapunk:
        \begin{align*}
          h'(x)=h(x)\bigg(&\frac{2x+3}{x^2+3x}
          +\coth x
          +\frac{1}{(1+x^2)\arctan x}\\
          &+\frac{\sin x}{1+\cos x}\bigg)
          \text.
        \end{align*}
        Ez az alak minden olyan intervallumon használható, ahol az eredeti
        függvény és a feltüntetett hányadosok értelmezettek. Megjegyzés: valós
        függvényként $\operatorname{arctanh}\pi$ nem értelmezett, ezért a
        feladat nevezőjében valószínűleg elírás van; a deriválás szempontjából
        mindenképpen konstansként viselkedik.

  \item Legyen $u(x)=\ln(\cos^2(x^4))$. Ekkor
        $$
          u'(x)=\frac{2\cos(x^4)(-\sin(x^4))4x^3}{\cos^2(x^4)}
          =-8x^3\tan(x^4)
          \text.
        $$
        Tehát
        $$
          \boxed{
          i'(x)=-\frac{8x^3\tan(x^4)}
          {3\,[\ln(\cos^2(x^4))]^{2/3}}}
          \text.
        $$

  \item Írjuk
        $$
          u(x)=\sqrt{\frac{x^2+3}{e^{2x}}}
          =e^{-x}\sqrt{x^2+3}
          \text.
        $$
        Logaritmikus deriválással
        $$
          \frac{u'}u=\frac{x}{x^2+3}-1
          \text.
        $$
        Mivel $(\ln(\arcsin u))'=u'/[\arcsin u\sqrt{1-u^2}]$,
        $$
          \boxed{
          j'(x)=\frac{u(x)\left(\dfrac{x}{x^2+3}-1\right)}
          {\arcsin(u(x))\sqrt{1-u(x)^2}}}
          \text,
        $$
        ahol a képlet az eredeti függvény értelmezési tartományának belső
        pontjain érvényes.

  \item $x>0$ esetén $x^x=e^{x\ln x}$, ezért
        $$
          \boxed{(x^x)'=x^x(\ln x+1)}
          \text.
        $$

  \item Mivel $(x^2+1)^{\sin x}
        =e^{\sin x\ln(x^2+1)}$,
        $$
          \boxed{
          l'(x)=(x^2+1)^{\sin x}
          \left(\cos x\ln(x^2+1)+\frac{2x\sin x}{x^2+1}\right)}
          \text.
        $$
\end{enumerate}

\subsection{Magasabb rendű deriváltak}

\begin{enumerate}[a)]
  \item Ha $m\in\mathbb N_0$, akkor
        $$
          \boxed{
          \frac{\mathrm d^n}{\mathrm dx^n}x^m=
          \begin{cases}
            \dfrac{m!}{(m-n)!}x^{m-n},&0\leq n\leq m,\\
            0,&n>m.
          \end{cases}}
          \text.
        $$
        Különösen az $m$-edik derivált $m!$, nem pedig $1$.

  \item A deriváltak négyes periódusban ismétlődnek:
        $$
          \sin x,\quad\cos x,\quad-\sin x,\quad-\cos x,\quad\sin x,\dots
        $$
        Ezt egyetlen képlet foglalja össze:
        $$
          \boxed{\displaystyle
          \frac{\mathrm d^n}{\mathrm dx^n}\sin x
          =\sin\left(x+\frac{n\pi}{2}\right)}
          \text.
        $$
\end{enumerate}

\subsection{Érintő és normális}

Legyen
$$
  f(x)=2x^3+3\sqrt x-\frac{3}{2x^2}
  \text.
$$
Az $x_0=1$ pontban
$$
  f(1)=2+3-\frac32=\frac72
  \text,
$$
és
$$
  f'(x)=6x^2+\frac{3}{2\sqrt x}+3x^{-3},
  \qquad
  f'(1)=6+\frac32+3=\frac{21}{2}
  \text.
$$
Az érintő egyenlete
$$
  y-\frac72=\frac{21}{2}(x-1),
  \qquad
  \boxed{y=\frac{21x-14}{2}}
  \text.
$$
A normális meredeksége a negatív reciprok, $-2/21$, ezért
$$
  y-\frac72=-\frac2{21}(x-1),
  \qquad
  \boxed{y=-\frac2{21}x+\frac{151}{42}}
  \text.
$$

\clearpage
\subsection{A kör megadott irányú érintői}

A $3x-4y+7=0$ egyenes meredeksége $3/4$. Az
$x^2+y^2=25$ kör implicit deriválásából
$$
  2x+2yy'=0,
  \qquad y'=-\frac{x}{y}
  \text.
$$
A párhuzamosság feltétele
$$
  -\frac{x}{y}=\frac34,
  \qquad y=-\frac43x
  \text.
$$
Ezt a kör egyenletébe helyettesítve
$$
  x^2+\frac{16}{9}x^2=25,
  \qquad x^2=9
  \text.
$$
Így a keresett pontok
$$
  \boxed{(3,-4)\quad\text{és}\quad(-3,4)}
  \text.
$$

\end{document}
